\section{Независимая проверка устойчивости к оформлению этапов}
\label{sec:segregation80}

Основной контраст SR--SN операционализирует ролевую организацию заменой обозначений. Дополнительный проспективный эксперимент проверяет устойчивость такого сравнения при другом, явно контролируемом оформлении инструкций. Все условия используют один новый ответ; ролевые обозначения (planner, implementer, reviewer либо stage 1, stage 2, stage 3) скрещиваются с заголовками в квадратных скобках на отдельных строках либо обозначениями с двоеточием в одном абзаце. Внутри каждого оформления меняются только три обозначения; между оформлениями сохраняются предложения с операциями и их порядок. Общая инструкция требует внутреннего анализа, реализации и проверки и запрещает комментарии об этапах в ответе. Каждый кандидат должен быть одной полной программой в блоке Python. Топология вызовов и компрессия не варьируются.

Такое воздействие изолирует замену ролевых обозначений при каждом оформлении, но не удостоверяет выполнение отдельных внутренних ролей: соблюдение этапов и приватные вычисления модели ненаблюдаемы. Поэтому сравнение оформлений относится к видимой структуре инструкции, а не к независимым агентам или раздельному доступу к информации. Сохранение эффекта обозначений при обоих оформлениях усилило бы вывод об устойчивости данной реализации; смена направления или широкий интервал ограничивали бы такой вывод. Дополнительные данные не включаются в исходное семейство четырёх тестов.

Выборка содержит следующие 80 идентификаторов в ранее рандомизированном резервном порядке после исключения всех 200 основных задач, 40 задач разработки и восьми раскрытых примеров карточки датасета. Механический отбор предшествует генерации. Каждая задача получает четыре условия: всего 320 назначений и не более 80 парных единиц. Закреплённый порядок блоков задач чередует условия при четырёх параллельных вызовах. Сохраняются GPT-5.4 mini, medium reasoning, Codex CLI 0.153.4, новые эфемерные вызовы через подписку, полный предоставленный контекст и отключённые инструменты. Лимит вызова составляет 600 секунд, ограничение входа --- 65\,536 байт; запросы не сокращаются, кандидаты не запускаются повторно. Лимит времени отличается от основной серии, поэтому дополнительные оценки приводятся отдельно.

Исходники штатного BigCodeBench и тесты не меняются. До генерации недостающие зависимости и старые библиотечные интерфейсы восстанавливаются в отдельно сохранённых образах и полных контрольных запусках. Итоговый образ проходит все 80 эталонных программ; неправильный контроль отклоняется на 79 задачах, но проходит на BigCodeBench/59. Эта задача остаётся назначенной и проверяется во всех условиях, однако качество по ней считается недоступным из-за контроля. Поэтому итоговый контроль допускает не более 79 пар для качества. Сохранены все неудачные контроли, версии пакетов и идентификаторы образов; эталонные программы, тесты и идентификаторы задач не заменяются. Процессы генерации не получают материалы оценщика.

Заранее заданное основное семейство содержит два точных двусторонних теста Мак-Немара: роли против нейтральных обозначений отдельно для заголовков и абзаца, с коррекцией Холма при семейном $\alpha=0.05$. Контрасты заголовков с абзацем и разность двух ролевых контрастов являются описательными. Бутстреп-интервалы используют 10\,000 повторных выборок задач и начальное состояние генератора 20260908. Ресурсы представлены парными разностями токенов, условной API-стоимости и отношениями парных средних; критерий эквивалентности, неухудшения или совместного выигрыша качества и стоимости не вводится. До генерации выполнен точный расчёт мощности при консервативном первом пороге Холма 0.025 для разных долей дискордантных пар и размеров эффекта. Выборка не обладает высокой мощностью для исключения малых эффектов; это ограничение сохраняется независимо от наблюдаемого результата. Полный план выполнения и анализа зафиксирован до запуска.
